\section{综合评价与适用范围}

\subsection{三问的收益与计算代价}
三种模型围绕同一决策接口逐步细化：问题一比较保留分量与释放内部并行的代价，问题二引入同核数据驻留，问题三展开共享读取的事件顺序。对应五核平均加速比分别为 3.8311、3.9522 和 4.1939。三者使用相同整图单核参照，平均值都先计算逐图比值再汇总；这些数字反映本题 100 图上的结果。

第一问的同场景完整分量对照用于观察扩展结构候选的收益，第二问与第一问的比较同时涉及运行机制和方案变化。第三问另行冻结同一方案切换 L2，能够直接测量该方案的缓存开关效果。因此，这三类比较回答的分别是候选扩展、跨场景方案表现和同方案硬件配置变化，分析时不能混作同一种消融实验。

算法以有限构造控制求解规模。图特征提取主要沿节点、边和张量关联表扫描；候选生成后的排序、联合图检查及驻留扫描随展开后的图规模增长。第三问进一步维护事件队列，使用更细的时间信息，400 次多核方案生成的进程内耗时中位数为 0.600 s、最大 10.916 s。该时间属于既有四工作进程实验条件，不包含官方验收和进程启动。

\subsection{独立输入与等价变换检验}
冷启动保证方案从当前输入重新构造；性能能否推广到其他输入，还需要额外测试。冻结后已构造串行链、独立分量、分支汇合、不规则分层、缓存压力及共享输入六类小图，并对部分图重编号和打乱数组次序。另加入 576、1152 个非 COPY 操作的两个分层图。测试图由固定种子生成，未从原 100 图抽取或拟合；正式评价在方案保存后进行。

独立测试揭示了第二问的稳定性不足。对同一个重编号后的不规则图，将原方案作等价 ID 映射得到 168070 cycle，重新求解得到 195740 cycle，耗时增加 16.46\%。两项方案在同一输入上验收，排除了仅由评估器面对不同编号所产生的直接差异；候选构造或选择本身也受编号与列表顺序影响。现有记录同时改变了这两项表示，尚不能进一步分离其贡献。

\begin{table}[htbp]
\centering
\caption{两个新分层图的五核独立测试；两问分别按自身场景验收}
\label{tab:newgraphs}
\begin{tabular}{rrrrr}
\hline
操作数 & 问题一/cycle & 问题二/cycle & 问题二/问题一 & 生成耗时一/二/s \\
\hline
576 & 773238 & 2318992 & 2.9991 & 0.4040 / 0.0312 \\
1152 & 1563459 & 4172607 & 2.6688 & 0.8077 / 0.0552 \\
\hline
\end{tabular}
\end{table}

两个新图中，第二问虽然生成更快，完成时间却明显增加。其五类候选少于第一问，驻留风险又使用固定系数，二者都可能影响选择；现有对照尚未识别各因素的独立贡献。结果说明场景 B 的复用条件并不自动转化为当前算法在任意输入上的优势。后续应针对候选覆盖和评分误差分开做消融，并使用新种子确认，不能把这两个反例用于调参后继续当作独立测试。

全部题面域内的已测正常输出均完成官方模拟。审计另含直接 Op--Op 边和非 COPY 搬运管道的越域样本，只用于程序兼容性压力检查，不计入题面域内的推广证据。问题三尚未完成同等级的新图检验，其全量结果与缓存机理验证不替代这一环节。

\subsection{机理近似与提交前核对}
第一问的工作量模型未显式展开动态换出，第二问的静态区间不能完整描述空间复用补边，第三问的事件图也未复制全部核内存储调度。因而三问共同的改进方向是校准“构造顺序—实际驻留—请求释放”之间的误差，而非继续扩大无差别枚举。对第二问尤其需要识别因编号打破并列而改变的方案，再评估结构特征能否提供更稳定的次序。

输入契约审计还发现，冻结的问题一、二入口没有前置拒绝重复 ID、负周期、重复边和仅经 COPY 闭合的非法环；自检也未完整覆盖人为注入的负数或布尔子图编号。实际构造器在已测正常输入上生成非负整数编号，原有验收成绩不受这些注入测试改变。面向未知来源图时仍应增加独立的字段、全图 DAG 和输出契约检查，该检查无需调用性能评估器。

方案文件、参数、输入摘要、生成记录及独立验收记录一并保留。论文中的图表由相应明细重新计算，逐图时间和额外搬运表随代码提供。已测范围内，有限构造能够在较短生成时间下取得多核收益；其边界集中在输入表示敏感性、驻留近似和未见结构的候选覆盖，需要以具体反例继续改进。
